% !TEX root =main.tex
\usepackage{nag}
\usepackage{graphicx}
\usepackage{amsthm}
% \usepackage{tabulary} % Unused; can interfere with standard tabular in RevTeX
\usepackage{qcircuit}
\usepackage{mathtools}
\usepackage{bm}
\usepackage[ruled,vlined]{algorithm2e}
%\usepackage[table,xcdraw]{xcolor}
\usepackage{braket}
\usepackage{datetime}
\usepackage{stackrel}
% Used to declare symbols below
% \usepackage{stmaryrd}
\usepackage{dsfont}
\usepackage{qcircuit}
\usepackage[english]{babel}
\usepackage{units}

\usepackage{tikz}
\usetikzlibrary{backgrounds,decorations.pathreplacing}

\usepackage{chngcntr}
\counterwithout{equation}{section}

\usepackage[normalem]{ulem}
\usepackage{slashed}
%%
\usepackage{soul}

\usepackage{xcolor}
% \usepackage[
%     colorlinks=true,
%     urlcolor=blue,
%     linkcolor=blue,
%     citecolor=blue,
%     filecolor=blue,
% ]{hyperref}
% \usepackage[caption=false]{subfig}
%%
\usepackage{amssymb}
\usepackage{amsmath}
\usepackage{amsthm}
\usepackage{braket}
\usepackage{mathdots}

% \usepackage[OT2,T1]{fontenc}
% \DeclareSymbolFont{cyrletters}{OT2}{wncyr}{m}{n}
% \DeclareMathSymbol{\Sha}{\mathalpha}{cyrletters}{"58}






\usepackage{makeidx}
\makeindex


% 

\usepackage{soul}
\usepackage{xargs}
%\usepackage{todonotes}
\newcommandx{\cmnote}[2][1=]{\linespread{1.0}\todo[linecolor=red,backgroundcolor=red!25,bordercolor=red,#1]{#2}}

\newcommand{\CS}{\text{CS}}
\newcommand{\vac}{\text{vac}}

% 
\let\underline\ul

% 

\allowdisplaybreaks[4]


\DeclareMathOperator{\tr}{tr}
\DeclareMathOperator{\Tr}{Tr}
\DeclareMathOperator{\TR}{TR}
\DeclareMathOperator{\Sp}{Sp}
\DeclareMathOperator{\spur}{sp}
\DeclareMathOperator{\diag}{diag}
\DeclareMathOperator{\Herm}{Herm}
\DeclareMathOperator{\cov}{cov}
\DeclareMathOperator{\var}{var}
\DeclareMathOperator{\Prob}{Prob}
\DeclareMathOperator{\spec}{spec}
\DeclareMathOperator{\sinc}{sinc}
\renewcommand{\Re}{\operatorname{Re}}
\renewcommand{\Im}{\operatorname{Im}}

\DeclareMathOperator{\sech}{sech}
\DeclareMathOperator{\csch}{csch}

\newcommand{\diffd}{\mathrm{d}}



 \index{} \index{} \index{} \index{} \index{} \index{} \index{} \index{}%%
\makeatletter

%

%
\newcommand{\ringplus}{\mathbin{\text{\@ringplus}}}

\newcommand{\@ringplus}{%
  \ooalign{\hidewidth\raise1.3ex\hbox{\tiny$\circ$}\hidewidth\cr$\m@th+$\cr}%
}

\newcommand{\ringminus}{\mathbin{\text{\@ringminus}}}

\newcommand{\@ringminus}{%
  \ooalign{\hidewidth\raise0.9ex\hbox{\tiny$\circ$}\hidewidth\cr$\m@th-$\cr}%
}
\makeatother
 \index{} \index{} \index{} \index{} \index{} \index{} \index{} \index{}%%

\newcommand{\rplus}{{\ringplus}}
\newcommand{\rminus}{{\ringminus}}

\newcommand{\nasc}{\epsilon}
\newcommand{\dnasc}{\delta_\nasc}
\newcommand{\tdnasc}{\tilde \delta_\nasc}

\newcommand{\qnasc}{q_\nasc}
\newcommand{\pnasc}{p_\nasc}


\newcommand{\mc}[1]{\mathcal{#1}}
\newcommand{\const}[0]{\text{(const.)}}
\newcommand{\cc}[0]{\text{c.c.}}
\newcommand{\bvec}[1]{\bm{#1}}

\newcommand{\tp}[0]{T}
\newcommand{\herm}[0]{\mathrm{H}}
\newcommand{\psinv}[0]{+}

\newcommand{\group}[1]{\mathrm{{#1}}}
\newcommand{\algebra}[1]{\mathfrak{{#1}}}
\newcommand{\logical}[0]{L}

\newcommand{\mathlatt}[1]{{\mathsf{#1}}}
\newcommand{\mathreclatt}[1]{\tilde{\mathsf{#1}}}
\newcommand{\dirlatt}{{\mathlatt L}}
\newcommand{\reclatt}{{\tilde \dirlatt}}
\newcommand{\vlatt}{{\mathlatt V}}
\newcommand{\flatt}{{\mathlatt F}}
\newcommand{\slatt}{{\mathlatt S}}


\newcommand{\weyl}{V}

\newcommand{\EPR}{\text{EPR}}
\newcommand{\bounceEPR}[2]{\op{A}_{#2}(#1)}
\newcommand{\bounceEPRgate}[2]{A_{#2}(#1)}

\newcommand{\dampop}[1]{\op{N}(#1)}
\newcommand{\dampgate}{N(\beta)}

% GKP commands
\newcommand{\GKP}{\text{GKP}} 
\newcommand{\GKPproj}{\op{\Pi}_\GKP}
\newcommand{\approxGKPproj}{\hat{\overline{\Pi}}_\GKP}
\newcommand{\qunaught}{\varnothing}
\newcommand{\sqz}{\grn \text{sqz}\blk}


\newcommand{\mvec}[1]{\tilde{\vec{#1}}}

\newcommand{\latG}{V}
%---------------------------------------------------------------------------
% Author commenting
%---------------------------------------------------------------------------




%---------------------------------------------------------------------------
% Sha character (Dirac comb)
%---------------------------------------------------------------------------

\DeclareFontFamily{U}{wncy}{}
\DeclareFontShape{U}{wncy}{m}{n}{<->wncyr10}{}
\DeclareSymbolFont{mcy}{U}{wncy}{m}{n}
\DeclareMathSymbol{\Sh}{\mathord}{mcy}{"58}



%---------------------------------------------------------------------------
% Specialized bra and ket defs
%---------------------------------------------------------------------------

% NOTE: basic bras and kets already defined in braket

% negative spacing for subscripts -- decide on one or two \! characters:
\newcommand{\negspace}{\!}
%\newcommand{\negspace}[0]{\!\!}

% Base commands for putting left and right subscripts on bras and kets and things
\newcommand{\lsub}[2]{{\protect\vphantom{#1}}_{#2} \negspace {#1}}
\newcommand{\rsub}[2]{{#1} \negspace {\protect\vphantom{#1}}_{#2}}
\newcommand{\lrsub}[3]{{\protect\vphantom{#1}}_{#2} \negspace {#1} \negspace {\protect\vphantom{#1}}_{#3}}

% Subscripted bras/kets
\newcommand{\ketsub}[2]{\rsub {\ket{#1}} {#2}}
\newcommand{\brasub}[2]{\lsub {\bra{#1}} {#2}}

% bra/ket macros
\newcommand{\pbra}[1]{\brasub{#1} p}
\newcommand{\qbra}[1]{\brasub{#1} q}
\newcommand{\pket}[1]{\ketsub{#1} p}
\newcommand{\qket}[1]{\ketsub{#1} q}



% basic inner product
\newcommand{\inprod}[2]{\left\langle {#1} | {#2} \right\rangle}
%\renewcommand{\braket}[0]{\inprod}
\newcommand{\mathinprod}[2]{\avg{{#1},{#2}}}

% doubly subscripted inner product
\newcommand{\inprodsubsub}[4]{\lrsub {\inprod{#1}{#2}} {#3} {#4}}

% inner product macros
\newcommand{\pqinprod}[2]{\inprodsubsub{#1}{#2} p q}
\newcommand{\qpinprod}[2]{\inprodsubsub{#1}{#2} q p}
\newcommand{\pinprod}[2]{\inprodsubsub{#1}{#2} p p}
\newcommand{\qinprod}[2]{\inprodsubsub{#1}{#2} q q}

%basic outer product
\newcommand{\outprod}[2]{\ket {#1}\!\bra {#2}}

% doubly subscripted outer product
\newcommand{\outprodsubsub}[4]{\ketsub {#1}{#3} \brasub{#2}{#4}}

% outer product macros
\newcommand{\pqoutprod}[2]{\outprodsubsub{#1}{#2}p q}
\newcommand{\qpoutprod}[2]{\outprodsubsub{#1}{#2}q p}
\newcommand{\qoutprod}[2]{\outprodsubsub{#1}{#2}q q}
\newcommand{\poutprod}[2]{\outprodsubsub{#1}{#2}p p}


%---------------------------------------------------------------------------
% Nascent-delta-based bra and ket defs
%---------------------------------------------------------------------------


% bra/ket macros
\newcommand{\pnascbra}[1]{\brasub{#1} \pnasc}
\newcommand{\qnascbra}[1]{\brasub{#1} \qnasc}
\newcommand{\pnascket}[1]{\ketsub{#1} \pnasc}
\newcommand{\qnascket}[1]{\ketsub{#1} \qnasc}

% inner product macros
\newcommand{\pqnascinprod}[2]{\inprodsubsub{#1}{#2} \pnasc \qnasc}
\newcommand{\qpnascinprod}[2]{\inprodsubsub{#1}{#2} \qnasc \pnasc}
\newcommand{\pnascinprod}[2]{\inprodsubsub{#1}{#2} \pnasc \pnasc}
\newcommand{\qnascinprod}[2]{\inprodsubsub{#1}{#2} \qnasc \qnasc}

% outer product macros
\newcommand{\pqnascoutprod}[2]{\outprodsubsub{#1}{#2}\pnasc \qnasc}
\newcommand{\qpnascoutprod}[2]{\outprodsubsub{#1}{#2}\qnasc \pnasc}
\newcommand{\qnascoutprod}[2]{\outprodsubsub{#1}{#2}\qnasc \qnasc}
\newcommand{\pnascoutprod}[2]{\outprodsubsub{#1}{#2}\pnasc \pnasc}




%---------------------------------------------------------------------------
% end of subscripted bra/ket defs
%---------------------------------------------------------------------------

\newcommand{\opbra}[1]{\left( {#1} \right\rvert}
\newcommand{\opket}[1]{\left\lvert {#1} \right)}
\newcommand{\opinprod}[2]{\left( {#1} \middle\vert {#2} \right)}
\newcommand{\opoutprod}[2]{\opket {#1} \opbra {#2}}



\newcommand{\avg}[1]{\left\langle {#1} \right\rangle}
\newcommand{\avgg}[1]{\langle {#1} \rangle}
\newcommand{\abs}[1]{\left\lvert{#1}\right\rvert}
\newcommand{\abss}[1]{\lvert{#1}\rvert}
\newcommand{\norm}[1]{\left\lVert{#1}\right\rVert}
\newcommand{\normm}[1]{\lVert{#1}\rVert}



\newcommand{\fpart}[1]{\left\{{#1}\right\}}
\newcommand{\fpartt}[1]{\{{#1}\}}

%\newcommand{\ipart}[1]{\left\lfloor{#1}\right\rfloor}
%\newcommand{\ipartt}[1]{\lfloor{#1}\rfloor}

\newcommand{\ipart}[1]{\left\lfloor{#1}\right\rceil}
\newcommand{\ipartt}[1]{\lfloor{#1}\rceil}


\newcommand{\cent}{C}
\newcommand{\sat}{S}



\DeclareMathOperator{\pv}{p.v.}




\newcommand{\reals}[0]{\mathbb{R}}
\newcommand{\complex}[0]{\mathbb{C}}
\newcommand{\wholenums}[0]{\mathbb{N}_0}
\newcommand{\integers}[0]{\mathbb{Z}}
\newcommand{\rationals}[0]{\mathbb{Q}}
\newcommand{\projective}{\mathbb{P}}


\newcommand{\op}[1]{\hat{#1}}
\newcommand{\ope}[1]{\hat{#1}}
%\newcommand{\op}[1]{#1}
\newcommand{\opvec}[1]{\op{\vec{#1}}}
\newcommand{\opmat}[1]{\op{\mat{#1}}}
%\newcommand{\id}[0]{1}
%\newcommand{\id}[0]{I}
\newcommand{\supop}[1]{\mathcal{#1}}

%\newcommand{\1}[0]{\mathit{1}}

%\newcommand{\1}[0]{\mat{\id}}

%\newcommand{\mat}[1]{\mathsf{#1}}
%\newcommand{\mat}[1]{\vec{#1}}
\newcommand{\mat}[1]{\bm{\mathrm{#1}}}
\newcommand{\barmat}[1]{\bar{\mat{#1}}}
\newcommand{\tildemat}[1]{\tilde{\mat{#1}}}
%\newcommand{\mat}[1]{{#1}}
\newcommand{\barvec}[1]{\bar{\vec{#1}}}
\newcommand{\tildevec}[1]{\tilde{\vec{#1}}}

%\renewcommand{\vec}[1]{\mathrm{#1}}
\renewcommand{\vec}[1]{\bm{#1}}

\newcommand{\controlled}[1]{\op{\mathrm{C}}_{#1}}
\newcommand{\CZ}[0]{\controlled Z}
\newcommand{\CX}[0]{\controlled X}
\newcommand{\CZct}[0]{\controlled Z^\rightarrow}
\newcommand{\CXct}[0]{\controlled X^\rightarrow}
\newcommand{\CZtc}[0]{\controlled Z^\leftarrow}
\newcommand{\CXtc}[0]{\controlled X^\leftarrow}

\newcommand{\opCZ}[0]{\op{\mathrm{C}}_{Z}}

\newcommand{\cH}[0]{\mathcal{H}}






 \index{} \index{} \index{} \index{} \index{} \index{} \index{} \index{}%%
% COLOUR COMMANDS
 \index{} \index{} \index{} \index{} \index{} \index{} \index{} \index{}%%
\newcommand{\grn}{\color{green!60!black}}
\newcommand{\red}{\color{red!90!black}}
\newcommand{\blu}{\color{blue!70!cyan}}
\newcommand{\blk}{\color{black}}
\newcommand{\gra}{\color{gray!50!white}}
\newcommand{\prp}{\color{purple!70!black}}
\newcommand{\org}{\color{orange!80!black}}
\newcommand{\pnk}{\color{pink}}

 \index{} \index{} \index{} \index{} \index{} \index{} \index{} \index{}%%
% VERY SPECIFIC COMMANDS
 \index{} \index{} \index{} \index{} \index{} \index{} \index{} \index{}%%
\usepackage{graphicx}
%\graphicspath{{./figures/}}

%Circuit macros
\newcommand{\bs}{\! B_{_{\! \downarrow}} \!\!}
%\newcommand{\bsop}{\op{B}_{_{\! \downarrow}}}
\newcommand{\bsop}{\op{B} }
% Figures Location

\newcommand{\rket}[1][0]{\ketsub{#1}{\theta}}
\newcommand{\rbra}[1][0]{\brasub{#1}{\theta}}
\usepackage{graphicx}
\usepackage{multirow}
\usepackage{hhline}
%\graphicspath{{../Figures/}}



%%%%%%%%%%%%%%%%%%%%%%%%%%%%%%%
% beamsplitter as arrow
% Use \bsbal{} just like \ctrl{}
%%%%%%%%%%%%%%%%%%%%%%%%%%%%%%%
\xyoption{color}
\xyoption{line}

\newcommand{\qwax}[1][-1]{\ar @{->} [#1,0]}
%Using the bsbal command requires setting the length. Specifying an optional argument before the length allows to change the colour, specifying an optional argument AFTER the length allows to change the line style (eg. .>, -->)
\newcommandx*\bsbal[3][1=black, 3=->]{\ar @[#1]@{#3} [#2,0] \qw}

\newcommandx*\varbs[5][1=black, 3=\theta,4=0.5,5=->]{\ar @[#1]@{#5}^(#4){#3} [#2,0] \qw}

%\newcommand{\bsbal}[1]{\qwax[{#1}] \qw}
\newcommand{\qwaxred}[1][-1]{\ar @[red]@{->} [#1,0]}
%\newcommand{\bsbalred}[1]{\qwaxred[{#1}] \qw}
\newcommand{\qwred}[1][-1]{\ar @[red]@{-} [0,#1]}
\newcommand{\gatered}[1]{*+<.6em>{#1} \POS ="i","i"+UR;"i"+UL **\dir{-};"i"+DL **\dir{-};"i"+DR **\dir{-};"i"+UR **\dir{-},"i" \qwred}


\newcommand{\ctrlnowire}[1]{\control \qwx[#1]}
\newcommandx*\lblline[3][3=0.5]{\ar @{-}^(#3){#1} [#2,0]}
% \newcommandx*\ctrlg[2]{\control \ar @{-}^{#1} [#2,0] \qw}
\newcommandx*\ctrlg[3][3=0.5]{ \raisebox{-3pt}{$\bullet$}  \ar @{-}^(#3){#1} [#2,0] \qw }
\newcommandx*\ctrlog[2]{\controlo \ar @{-}^{#1} [#2,0] \qw}
%%%Controlo gate followed by dashed wire
\newcommandx*\ctrlodash[1]{\controlo \ar @{-} [#1,0] \ar @[black]@{.} [0,-1]}
\newcommand{\qwdot}{\ar @[black]@{.} [0,-1]}
%%%%%%%%%%%%%%%%%%%%%%%%%%%%%%%

%%%%%%%%%%%%%%%%%%%%%%%%%%%%%%%
%EPR state angled wire
%use \EPR followed by u, d for up and down and then r, l for right and left
%%%%%%%%%%%%%%%%%%%%%%%%%%%%%%%

\newcommand{\EPRdr}{\ar @{-} [dr(0.5)]\qw }
\newcommand{\EPRur}{\ar @{-} [ur(0.5)]\qw }
\newcommand{\EPRul}{\ar @{-} [ul(0.5)] }
\newcommand{\EPRdl}{\ar @{-} [dl(0.5)] }

%%%%%%%%%% invisible gate box %%%%%%%%%%%%%
\newcommand{\nogate}[1]{*+<.6em>{#1} \POS ="i","i"+UR;"i"+UL **[white]\dir{-};"i"+DL **[white]\dir{-};"i"+DR **[white]\dir{-};"i"+UR **[white]\dir{-},"i" }

%%%%%%%%%%%%% Multigate invisible, use \nghost for lower sections, not \ghost %%%%%%%%%%%%%%
\newcommand{\nomultigate}[2]{*+<1em,.9em>{\hphantom{#2}} \POS [0,0]="i",[0,0].[#1,0]="e",!C *{#2},"e"+UR;"e"+UL **[white]\dir{-};"e"+DL **[white]\dir{-};"e"+DR **[white]\dir{-};"e"+UR **[white]\dir{-},"i"}

\newcommand{\colorgategroup}[7]{\POS"#1,#2"."#3,#2"."#1,#4"."#3,#4"!C*+<#5>[#7]\frm{#6}}

\setcounter{page}{1}
\pagenumbering{arabic}

\newcommand{\prlsection}[1]{\vspace{0.5ex}\textbf{#1}}

\newcommand*\circled[1]{\tikz[baseline=(char.base)]{
            \node[shape=circle,draw,inner sep=2pt] (char) {#1};}}
